% =============================================================================
% Sección: Implementación distribuida con MPI
% =============================================================================

\section{Desarrollo de la implementación}

Esta sección describe cómo se construyó el simulador de tráfico unidimensional
con paralelismo MPI, explicando las decisiones de diseño y sus consecuencias en
rendimiento. La idea central fue mantener separadas la lógica del autómata, la
gestión del estado distribuido y el mecanismo de comunicación, de modo que
cualquier diferencia de rendimiento entre configuraciones sea atribuible al
número de procesos o a la estrategia de compilación, y no a cambios en el
algoritmo.

% ---------------------------------------------------------------------------
\subsection{Modelo del autómata celular de tráfico}
% ---------------------------------------------------------------------------

El sistema modela una carretera circular de $N$ celdas. Cada celda
$i \in \{0, 1, \ldots, N-1\}$ puede estar vacía ($R = 0$) o contener un carro
($R = 1$). En cada paso de tiempo, un carro avanza a la celda siguiente si está
libre, o se queda donde está si está ocupada. La regla de actualización es:

\begin{equation}
    R(t+1,\, i) =
    \bigl[R(t,i) \wedge R(t,\,(i+1)\bmod N)\bigr]
    \;\vee\;
    \bigl[R(t,\,(i-1)\bmod N) \wedge \neg R(t,i)\bigr].
    \label{eq:regla}
\end{equation}

El primer término representa un carro que no puede avanzar porque la celda
siguiente está ocupada. El segundo representa un carro de la celda anterior que
avanza hacia una celda libre. La condición $\neg R(t,i)$ en el segundo término
evita que dos carros ocupen la misma celda simultáneamente.

Cada celda se almacena como un entero sin signo de 8 bits (\texttt{uint8\_t}).
La regla anterior se traduce a operaciones de bits:

\begin{equation}
    \texttt{next}[i] =
    (h \;\&\; a) \;\mid\;
    (b \;\&\; \mathtt{(uint8\_t)(\sim h)} \;\&\; \mathtt{1u}),
    \label{eq:kernel}
\end{equation}

donde $b = R(t,\,i{-}1)$, $h = R(t,\,i)$ y $a = R(t,\,i{+}1)$. El cast
explícito \texttt{(uint8\_t)(\textasciitilde h)} es necesario porque el operador
\texttt{\textasciitilde} en C promueve el operando a \texttt{int} antes de
invertirlo; sin el cast, el resultado sería \texttt{0xFFFFFFFE} en lugar de
\texttt{0xFE}, y la conjunción con \texttt{1u} produciría cero
incorrectamente cuando $h = 1$.

Las lecturas del paso $t$ provienen del arreglo \texttt{cells} y las escrituras
del paso $t+1$ van al arreglo \texttt{next\_cells}. Como son arreglos distintos
y cada posición nueva depende únicamente de valores del paso anterior, todas las
iteraciones del bucle son completamente independientes entre sí, lo que hace el
problema embarazosamente paralelo.

% ---------------------------------------------------------------------------
\subsection{Velocidad media y estado estacionario}
% ---------------------------------------------------------------------------

La velocidad media en el paso $t$ es la fracción de carros que se desplazaron:

\begin{equation}
    v(t) = \frac{1}{N_c}
    \sum_{i=0}^{N-1}
    \mathbf{1}\bigl[R(t,i) = 1 \;\wedge\; R(t+1,i) = 0\bigr],
    \label{eq:velocidad}
\end{equation}

donde $N_c$ es el número total de carros. Este valor es constante durante toda
la simulación: la regla~\eqref{eq:regla} nunca crea ni elimina carros, lo que
se verifica explícitamente en los tests de corrección.

Partiendo de una distribución inicial aleatoria, el sistema atraviesa una fase
transitoria donde $v(t)$ fluctúa mientras los carros se reorganizan. Tras
suficientes pasos, $v(t)$ converge a la velocidad asintótica $v_\infty$, que es
el valor de interés. El benchmark reporta el promedio sobre la fase de medición:

\begin{equation}
    v_\infty \approx \bar{v} =
    \frac{1}{T_m}
    \sum_{t = T_w + 1}^{T_w + T_m} v(t),
    \label{eq:vasint}
\end{equation}

donde $T_w$ es el número de pasos de calentamiento y $T_m$ el número de pasos
de medición.

% ---------------------------------------------------------------------------
\subsection{Criterio de convergencia del calentamiento}
% ---------------------------------------------------------------------------

La duración de la fase transitoria varía con la densidad y con $N$. Para
detectar la convergencia automáticamente, se compara la velocidad media en dos
ventanas consecutivas de $W$ pasos. El calentamiento termina cuando:

\begin{equation}
    \left|
        \bar{v}_{[t-W,\, t]}
        - \bar{v}_{[t-2W,\, t-W]}
    \right| < \varepsilon,
    \label{eq:criterio}
\end{equation}

con $W = 50$, $\varepsilon = 10^{-4}$ y un mínimo de 200 pasos antes de evaluar
el criterio, para evitar falsos positivos al inicio. El parámetro
\texttt{max\_warmup\_steps} actúa como techo; si la condición~\eqref{eq:criterio}
no se cumple antes de ese límite, el calentamiento termina de todas formas. En
el script de benchmark este techo se fija como $\max(N/10,\; 2000)$, que
garantiza al menos el tiempo suficiente para que una perturbación recorra la
carretera completa.

% ---------------------------------------------------------------------------
\subsection{Modelo de distribución de datos}
% ---------------------------------------------------------------------------

Con $P$ procesos MPI y una carretera de $N$ celdas, el proceso de rango $r$
recibe un bloque contiguo de celdas indexadas globalmente en
$[\mathit{offset}_r,\; \mathit{offset}_r + \ell_r - 1]$. La distribución se
calcula con una partición por bloques que reparte el residuo entre los primeros
rangos:

\begin{equation}
    \ell_r = \left\lfloor \frac{N}{P} \right\rfloor
             + \mathbf{1}[r < N \bmod P],
    \qquad
    \mathit{offset}_r = r \left\lfloor \frac{N}{P} \right\rfloor
                        + \min(r,\; N \bmod P).
    \label{eq:distribucion}
\end{equation}

Esta distribución garantiza que la diferencia entre el bloque más largo y el más
corto sea como máximo una celda, manteniendo el balance de carga para el kernel
de cómputo, cuyo trabajo por celda es completamente uniforme. La implementación
calcula los arreglos de conteos y desplazamientos en la función
\texttt{road\_mpi\_build\_distribution}, que se usa tanto en el scatter inicial
como en el gather del validador.

El estado global se genera en el rango~0 usando
\texttt{srand((unsigned~int)time(NULL))} y se dispersa a todos los procesos con
\texttt{MPI\_Scatterv}. Para que la implementación sea correcta se requiere
$P \leq N$; la función \texttt{road\_mpi\_create} verifica esta condición y
aborta con un mensaje diagnóstico si se viola, ya que un bloque de longitud cero
rompe la topología en anillo descrita en la Sección~\ref{sec:frontera}.

% ---------------------------------------------------------------------------
\subsection{Estructura de datos distribuida}
% ---------------------------------------------------------------------------

La estructura \texttt{RoadMPI} encapsula el estado completo de cada proceso:

\begin{itemize}
    \item \texttt{cells} y \texttt{next\_cells}: búferes de \texttt{uint8\_t}
          reservados con \texttt{posix\_memalign(64,~$\cdot$)}, lo que alinea
          cada búfer al inicio de una línea de caché de 64 bytes. Esta
          alineación beneficia al prefetcher y permite al compilador emitir
          cargas AVX2 alineadas en el bucle de actualización.

    \item \texttt{global\_length}, \texttt{local\_length},
          \texttt{global\_offset}: parámetros de la
          distribución~\eqref{eq:distribucion}.

    \item \texttt{global\_car\_count}: número total de carros $N_c$, calculado
          una sola vez con \texttt{MPI\_Allreduce} tras la dispersión inicial y
          constante durante toda la simulación. Se usa como denominador
          en~\eqref{eq:velocidad}.

    \item \texttt{left\_rank} y \texttt{right\_rank}: rangos del vecino
          izquierdo $(r-1) \bmod P$ y derecho $(r+1) \bmod P$. La aritmética
          modular establece la topología en anillo que implementa las condiciones
          de frontera periódicas.

    \item \texttt{comm}: comunicador MPI, parametrizable para facilitar pruebas
          con subcomunicadores.
\end{itemize}

Al finalizar cada paso de simulación, los punteros \texttt{cells} y
\texttt{next\_cells} se intercambian directamente sin copiar datos, de la misma
forma que en la implementación secuencial. No se realizan asignaciones de memoria
dentro del bucle de simulación, lo que elimina ese ruido de las mediciones de
tiempo.

% ---------------------------------------------------------------------------
\subsection{Condiciones de frontera periódicas en el entorno distribuido}
\label{sec:frontera}
% ---------------------------------------------------------------------------

La carretera es circular: la celda~$0$ tiene como vecina izquierda a la
celda~$N-1$, y la celda~$N-1$ tiene como vecina derecha a la celda~$0$. En el
entorno distribuido, la celda~$0$ local del proceso~$r$ y la celda~$\ell_r - 1$
local del proceso~$r$ necesitan información de procesos distintos para que la
regla~\eqref{eq:kernel} pueda aplicarse:

\begin{itemize}
    \item El proceso $r$ necesita la última celda del proceso $r-1$ como
          \emph{halo izquierdo} al calcular $R^{t+1}(\mathit{offset}_r)$.
    \item El proceso $r$ necesita la primera celda del proceso $r+1$ como
          \emph{halo derecho} al calcular
          $R^{t+1}(\mathit{offset}_r + \ell_r - 1)$.
\end{itemize}

Para $r = 0$, el halo izquierdo proviene del proceso $P-1$ (borde derecho de
la carretera global), y para $r = P-1$, el halo derecho proviene del
proceso~$0$ (borde izquierdo). Esto implementa exactamente la condición
periódica $R^t(0) \equiv R^t(N)$ de forma transparente al kernel de cómputo.

% ---------------------------------------------------------------------------
\subsection{Intercambio de halos no bloqueante}
% ---------------------------------------------------------------------------

Antes de cada paso de simulación, cada proceso envía sus celdas de borde a sus
vecinos y recibe las de ellos. La función \texttt{post\_halo\_exchange} publica
cuatro operaciones no bloqueantes con un esquema de etiquetas que evita
ambigüedades:

\begin{itemize}
    \item \texttt{TAG\_SEND\_LEFT} (100): etiqueta con la que un proceso envía
          \texttt{cells[0]} a su vecino izquierdo (\texttt{left\_rank}), para
          que este lo use como halo derecho.
    \item \texttt{TAG\_SEND\_RIGHT} (101): etiqueta con la que un proceso envía
          \texttt{cells[$\ell_r - 1$]} a su vecino derecho (\texttt{right\_rank}),
          para que este lo use como halo izquierdo.
\end{itemize}

La simetría del esquema establece que el proceso $r$ recibe su halo izquierdo
del proceso $r-1$ con \texttt{TAG\_SEND\_RIGHT} (porque $r-1$ está enviando su
celda más a la derecha), y recibe su halo derecho del proceso $r+1$ con
\texttt{TAG\_SEND\_LEFT} (porque $r+1$ está enviando su celda más a la
izquierda). La función publica las cuatro operaciones y retorna los cuatro
objetos \texttt{MPI\_Request} al llamador:

\begin{center}
\begin{verbatim}
MPI_Irecv(left_halo,  1, MPI_UINT8_T, left_rank,  TAG_SEND_RIGHT, ...);
MPI_Irecv(right_halo, 1, MPI_UINT8_T, right_rank, TAG_SEND_LEFT,  ...);
MPI_Isend(&cells[0],       1, MPI_UINT8_T, left_rank,  TAG_SEND_LEFT,  ...);
MPI_Isend(&cells[len - 1], 1, MPI_UINT8_T, right_rank, TAG_SEND_RIGHT, ...);
\end{verbatim}
\end{center}

El uso de operaciones no bloqueantes (\texttt{Irecv}/\texttt{Isend}) en lugar de
sus equivalentes bloqueantes es deliberado: permite solapar la latencia de red
con el cómputo del interior del bloque local, como se describe en la
sección siguiente.

% ---------------------------------------------------------------------------
\subsection{Kernel distribuido con solapamiento comunicación-cómputo}
% ---------------------------------------------------------------------------

La función \texttt{compute\_next\_and\_count\_distributed} ejecuta un paso
completo del autómata en cinco fases:

\begin{enumerate}
    \item \textbf{Publicar halos.} Se llama a \texttt{post\_halo\_exchange}.
          Las cuatro operaciones quedan en vuelo mientras el proceso
          continúa trabajando localmente.

    \item \textbf{Cómputo del interior.} Se aplica la regla~\eqref{eq:kernel}
          a las celdas $i \in [1,\, \ell_r - 2]$, que solo dependen de vecinos
          locales. Para $\ell_r \geq 3$ este bucle se ejecuta sin necesidad de
          esperar la red:

\begin{center}
\begin{verbatim}
for (int i = 1; i < len - 1; i++) {
    nxt[i]     = update_cell(cur[i-1], cur[i], cur[i+1]);
    departed  += (cur[i] == 1u && nxt[i] == 0u) ? 1 : 0;
}
\end{verbatim}
\end{center}

          El conteo de carros desplazados se acumula en la misma pasada,
          evitando un segundo recorrido sobre los datos. Los casos $\ell_r = 1$
          y $\ell_r = 2$ se tratan después de la espera de halos, ya que en
          esos casos todas las celdas son bordes.

    \item \textbf{Esperar halos.} \texttt{MPI\_Waitall} bloquea hasta que las
          cuatro operaciones publicadas en la fase~1 se completan. Si la red
          entregó los datos durante la fase~2, el tiempo de espera efectivo
          es cero.

    \item \textbf{Cómputo de los bordes.} Se aplica la regla~\eqref{eq:kernel}
          a las celdas $i = 0$ e $i = \ell_r - 1$ usando los halos recibidos.

    \item \textbf{Reducción global.} \texttt{MPI\_Allreduce} suma los contadores
          locales de carros desplazados y entrega el total a todos los procesos:
          \[
              v(t) = \frac{\displaystyle\sum_{r=0}^{P-1}
              \mathit{departed}_r}{N_c}.
          \]
          Todos los procesos reciben exactamente el mismo valor de $v(t)$,
          lo que es fundamental para el criterio de convergencia del
          calentamiento.
\end{enumerate}

\begin{figure}[H]
\centering
\fbox{%
\begin{minipage}{0.94\textwidth}
\centering
\textbf{Flujo de un paso distribuido (proceso $r$)}\\[4pt]
\texttt{Irecv}/\texttt{Isend} halos
$\;\longrightarrow\;$ cómputo interior $[1,\, \ell_r{-}2]$
$\;\longrightarrow\;$ \texttt{Waitall}
$\;\longrightarrow\;$ cómputo bordes $\{0,\; \ell_r{-}1\}$
$\;\longrightarrow\;$ \texttt{Allreduce}($\mathit{departed}$)
$\;\longrightarrow\;$ intercambio de punteros
\end{minipage}}
\caption{Ciclo de un paso de simulación en la implementación MPI.}
\label{fig:flujo_mpi}
\end{figure}

% ---------------------------------------------------------------------------
\subsection{Propiedad SPMD del calentamiento}
% ---------------------------------------------------------------------------

Como todos los procesos obtienen exactamente el mismo valor de $v(t)$ a través
de \texttt{MPI\_Allreduce}, la función
\texttt{simulation\_mpi\_warmup\_until\_steady\_state} evalúa el
criterio~\eqref{eq:criterio} de forma independiente en cada rango y llega a la
misma conclusión en el mismo paso, sin necesitar ninguna comunicación adicional
para sincronizar el fin del calentamiento. Esta propiedad, que se denomina
\emph{SPMD-idéntica}, simplifica el código del calentamiento y garantiza que
\texttt{MPI\_Finalize} se llame colectivamente sin interbloqueo en cualquier
configuración.

% ---------------------------------------------------------------------------
\subsection{Temporización y salida}
% ---------------------------------------------------------------------------

La medición de tiempo usa \texttt{MPI\_Wtime}, que en OpenMPI resuelve a un
contador de hardware de alta resolución sincronizado entre los procesos del mismo
trabajo. El cronómetro se activa justo antes de la fase de medición, dejando el
calentamiento fuera del tiempo reportado. Como los procesos no terminan
exactamente al mismo tiempo, se aplica \texttt{MPI\_Reduce} con \texttt{MPI\_MAX}
para reportar el tiempo del proceso más lento:

\begin{equation}
    t_{\text{wall}} = \max_{r \in [0,P)}
    \bigl(t_{\text{fin},r} - t_{\text{inicio},r}\bigr).
    \label{eq:tiempo_mpi}
\end{equation}

Este valor refleja la duración real de la ejecución sincronizada y es el
apropiado para calcular el speedup. El rango~0 imprime exactamente una línea
en \texttt{stdout}:

\begin{center}
\texttt{TIME\_MS AVG\_VELOCITY}
\end{center}

Los mensajes diagnósticos (número de pasos de calentamiento, número de procesos
efectivos) se envían a \texttt{stderr} para no interferir con la captura de
resultados del script de benchmark.

La interfaz de línea de comandos es:

\begin{center}
\texttt{traffic\_mpi\_opt} $N$ $\rho$ \texttt{max\_warmup} $T_m$
\end{center}

% ---------------------------------------------------------------------------
\subsection{Arquitectura de módulos}
% ---------------------------------------------------------------------------

La implementación se divide en cuatro módulos con responsabilidades bien
delimitadas:

\begin{table}[H]
\centering
\caption{Módulos de la implementación MPI y sus responsabilidades.}
\label{tab:modulos_mpi}
\begin{tabular}{lll}
\hline
Archivo & Cabecera & Responsabilidad \\
\hline
\texttt{road\_mpi.c}       & \texttt{road\_mpi.h}       & Estado distribuido, scatter, gather \\
\texttt{simulation\_mpi.c} & \texttt{simulation\_mpi.h} & Kernel temporal, halos, Allreduce \\
\texttt{mpi\_main.c}       & ---                        & Orquestación, CLI, temporización \\
\texttt{validator\_mpi.c}  & \texttt{validator\_mpi.h}  & Tests de corrección distribuidos \\
\hline
\end{tabular}
\end{table}

El módulo \texttt{mpi\_common.h} define las etiquetas de mensajes y los
helpers \texttt{left\_rank}/\texttt{right\_rank} como funciones estáticas en
línea, compartidas entre \texttt{road\_mpi.c} y \texttt{simulation\_mpi.c}.

% ---------------------------------------------------------------------------
\subsection{Variantes de compilación}
% ---------------------------------------------------------------------------

Se generan dos binarios que difieren únicamente en las flags de compilación.
El compilador base es \texttt{mpicc}, que es un wrapper sobre \texttt{gcc} que
inyecta automáticamente los flags de include y linker para la implementación MPI
disponible; todos los flags de optimización llegan directamente al gcc
subyacente.

\begin{table}[H]
\centering
\caption{Binarios generados por \texttt{make mpi mpi\_opt}.}
\label{tab:binarios_mpi}
\begin{tabular}{llcl}
\hline
Binario & \texttt{impl} en CSV & Flags OPT & Descripción \\
\hline
\texttt{traffic\_mpi}      & \texttt{traffic\_mpi}     & No  & Referencia sin optimización \\
\texttt{traffic\_mpi\_opt} & \texttt{traffic\_mpi\_opt} & Sí & Producción: \texttt{-O3 -march=native -flto} \\
\hline
\end{tabular}
\end{table}

Las flags de optimización actúan sobre el mismo perfil de kernel que en la
versión secuencial: tres cargas de memoria, dos operaciones lógicas y una
escritura por celda, limitado por el ancho de banda de memoria para $N$ grande.

\begin{itemize}
    \item \texttt{-O3}: activa vectorización automática del bucle interior,
          transformaciones de bucle e inlining agresivo.

    \item \texttt{-march=native}: emite instrucciones AVX2 o AVX-512 según el
          procesador. Con datos \texttt{uint8\_t}, un registro AVX2 de 256 bits
          contiene 32 celdas, permitiendo procesar 32 posiciones por instrucción
          SIMD.

    \item \texttt{-flto} con \texttt{-fno-semantic-interposition}: permite
          inlinar \texttt{update\_cell} y el kernel completo directamente en el
          bucle principal, eliminando el overhead de llamada en cada uno de los
          $T_m = 1000$ pasos de medición. \texttt{mpicc} delega el linking a
          \texttt{gcc}, por lo que \texttt{-flto} funciona correctamente con este
          wrapper.

    \item \texttt{-funroll-loops}: reduce el overhead del contador del bucle;
          especialmente beneficioso para bloques locales pequeños que caben en
          L1.

    \item \texttt{-falign-loops=32}: alinea el inicio del bucle interior a un
          boundary de línea de instrucción AVX2, evitando que el cuerpo del
          bucle cruce dos líneas de caché de instrucciones.
\end{itemize}

% ---------------------------------------------------------------------------
\subsection{Validador de corrección distribuido}
% ---------------------------------------------------------------------------

El validador \texttt{traffic\_validator\_mpi} ejecuta siete tests que deben
pasar con cualquier número de procesos $P \in \{1, 2, 4\}$:

\begin{enumerate}
    \item \textbf{Carretera vacía.} Con densidad $\rho = 0$,
          \texttt{global\_car\_count} $= 0$ y la velocidad debe ser exactamente
          $0.0$.

    \item \textbf{Carretera llena.} Con densidad $\rho = 1$, cada carro está
          bloqueado y la velocidad debe ser exactamente $0.0$.

    \item \textbf{Un solo carro.} Con un único carro en una carretera libre, la
          velocidad debe ser exactamente $1.0$.

    \item \textbf{Frontera periódica.} Un carro en la última celda global debe
          aparecer en la celda~0 tras un paso, cruzando la frontera entre el
          proceso $P-1$ y el proceso~0.

    \item \textbf{Conservación del número de carros.} El conteo de celdas
          ocupadas en el estado reunido con \texttt{road\_mpi\_gather\_cells}
          debe ser constante durante 20 pasos.

    \item \textbf{Coincidencia con referencia calculada a mano.} Para el estado
          inicial $\{0, 1, 0, 0, 1, 0, 1, 0\}$ se calcula a mano el estado
          esperado tras un paso aplicando la regla~\eqref{eq:kernel} con
          condiciones periódicas y se compara contra el resultado distribuido.
          Este test verifica simultáneamente la regla, el intercambio de halos
          entre procesos y el gather.

    \item \textbf{Avance exacto de un carro solitario.} Un carro en la celda~0
          debe estar en la celda $t \bmod N$ tras $t$ pasos, verificando el
          wrap-around completo a lo largo de $N$ iteraciones.
\end{enumerate}

Los tests que requieren comparar el estado global usan
\texttt{road\_mpi\_gather\_cells} para reunir los datos en el rango~0, que
evalúa la condición y transmite el resultado a todos los rangos con
\texttt{MPI\_Bcast}. Esto garantiza que todos los procesos lleguen al mismo
veredicto y que \texttt{MPI\_Finalize} se llame colectivamente sin interbloqueo.

El estado inicial de cada test se inyecta con
\texttt{road\_mpi\_create\_with\_state}, que acepta un arreglo global en el
rango~0 y lo distribuye con el mismo \texttt{MPI\_Scatterv} que usa la
inicialización aleatoria, permitiendo probar condiciones exactas sin depender
del generador de números aleatorios.

% ---------------------------------------------------------------------------
\subsection{Despliegue en el clúster AWS}
% ---------------------------------------------------------------------------

Las pruebas se ejecutaron en un clúster de cuatro instancias EC2
\texttt{c5.xlarge}, cada una con 4~vCPUs y 8~GiB de RAM. Todas las instancias
pertenecen a la misma zona de disponibilidad y están conectadas dentro de la
misma VPC privada de AWS, lo que garantiza latencias de red inferiores a
$100\,\mu\mathrm{s}$ entre nodos y un ancho de banda de hasta 10~Gbps, más que
suficiente para los halos de un byte por proceso y paso que intercambia este
autómata.

\subsubsection{Directorio compartido por NFS}

El nodo principal monta un volumen EFS (\emph{Elastic File System}) y lo
exporta por NFS al resto de los nodos en la ruta \texttt{/shared}. Allí se
almacenan los binarios compilados, el \texttt{hostfile} y el directorio de
resultados. De esta forma todos los nodos ejecutan exactamente el mismo binario
sin necesidad de copiarlo manualmente, y los CSV de resultados se escriben
directamente desde el nodo principal sin transferencias posteriores.

\subsubsection{Compilación y configuración de OpenMPI}

La compilación se realiza en el nodo principal con \texttt{mpicc} del paquete
\texttt{openmpi-bin}:

\begin{center}
\texttt{make mpi mpi\_opt validator}
\end{center}

Los binarios resultantes se depositan en \texttt{/shared/bin}, visible por
todos los nodos. El archivo \texttt{/shared/hostfile} declara cuántos procesos
puede ejecutar cada nodo:

\begin{center}
\begin{verbatim}
node0 slots=4
node1 slots=4
node2 slots=4
node3 slots=4
\end{verbatim}
\end{center}

Con cuatro nodos de cuatro núcleos cada uno, el barrido de procesos puede
alcanzar hasta $P = 16$ sin sobresubscripción. Para las pruebas de hasta
$P = 12$ del experimento de escalado siempre hay al menos un núcleo libre por
nodo para el sistema operativo, reduciendo la interferencia en las mediciones.

\subsubsection{Ejecución distribuida}

Cada medición se lanza desde el nodo principal con:

\begin{center}
\texttt{mpirun -{}-hostfile /shared/hostfile -{}-bind-to core -{}-map-by core} \\
\texttt{\quad -np $P$ /shared/bin/traffic\_mpi\_opt $N$ $\rho$ $T_{w,\max}$ $T_m$}
\end{center}

La opción \texttt{--bind-to core --map-by core} hace que OpenMPI asigne cada
proceso a un núcleo físico distinto en orden round-robin entre nodos. Para
$P \leq 4$ todos los procesos quedan en el nodo principal, eliminando el tráfico
de red y mostrando la escalabilidad dentro de un único nodo. Para $P > 4$ los
procesos se distribuyen entre nodos y el costo del intercambio de halos se vuelve
medible en la latencia, aunque mínimo en el ancho de banda dado que cada proceso
envía y recibe exactamente un byte por paso.

\subsubsection{Topología de comunicación}

Con la distribución por bloques~\eqref{eq:distribucion}, cada proceso
intercambia halos únicamente con sus dos vecinos inmediatos en el anillo. El
patrón de comunicación es una cadena de mensajes punto a punto de un byte, no
un \emph{all-to-all}. En AWS, donde la latencia de red entre instancias en la
misma zona de disponibilidad es inferior a $50\,\mu\mathrm{s}$, el costo de
comunicación es despreciable comparado con el costo de cómputo del bloque local
para $N \geq 10^5$.

% ---------------------------------------------------------------------------
\subsection{Diseño del script de benchmark}
% ---------------------------------------------------------------------------

El script \texttt{bench\_traffic\_mpi.sh} actúa como harness de medición: su
única responsabilidad es invocar los binarios con los parámetros correctos,
capturar su salida y escribirla en un CSV. El análisis y la generación de
gráficas son responsabilidad de \texttt{report.py}.

\subsubsection{Convención de la columna \texttt{np}}

La columna \texttt{np} del CSV registra el número de procesos MPI usados. Por
convención, $\mathtt{np} = 0$ indica la ejecución del binario secuencial
\texttt{traffic\_seq\_mem\_opt} sin MPI, que sirve como línea de referencia
absoluta del rendimiento de un único proceso sin overhead de comunicación.
Para $\mathtt{np} \geq 1$ el binario MPI se lanza con \texttt{mpirun -np NP}.

\subsubsection{Techo de calentamiento adaptativo}

El techo de pasos de calentamiento que el harness pasa a cada binario se
calcula como:

\begin{equation}
    T_{w,\max} = \max\!\left(\frac{N}{10},\; 2000\right).
    \label{eq:techo}
\end{equation}

El factor $N/10$ garantiza que el techo sea proporcional al tiempo que tarda
una perturbación en recorrer la carretera completa. El piso de 2000 evita
techos demasiado bajos para valores pequeños de $N$. El binario puede terminar
antes si el criterio~\eqref{eq:criterio} se cumple, lo que ocurre en la
práctica para densidades extremas en pocas centenas de pasos.

\subsubsection{Matriz de configuraciones}

El harness define dos experimentos que cubren las preguntas de rendimiento
centrales del trabajo.

El \textbf{experimento de escalado} mide todas las combinaciones de los seis
valores de $N$ con los conteos de procesos $\{0, 1, 2, 4, 8, 12\}$ a densidad
fija $\rho = 0.50$. Para $\mathtt{np} = 0$ se ejecuta el binario secuencial;
para $\mathtt{np} \geq 1$ se ejecutan \texttt{traffic\_mpi} y
\texttt{traffic\_mpi\_opt}, lo que permite separar la ganancia del paralelismo
de la ganancia del compilador.

El \textbf{experimento de densidad} fija un $N$ grande fuera del barrido
principal (80 millones de celdas, 160~MB, claramente DRAM-bound) y barre cinco
densidades $\{0.10, 0.30, 0.50, 0.70, 0.90\}$ con $\mathtt{np} \in \{0, 1, 4,
8\}$. El objetivo es verificar que la velocidad asintótica $v_\infty$ es
consistente entre configuraciones y que la distribución entre procesos no
introduce sesgo en la curva velocidad-densidad.

\subsubsection{Formato del CSV de resultados}

El archivo \texttt{data.csv} tiene siete columnas:

\begin{center}
\begin{tabular}{ll}
\hline
Columna & Descripción \\
\hline
\texttt{impl}           & Binario usado (véase Tabla~\ref{tab:binarios_mpi}) \\
\texttt{np}             & Número de procesos MPI (0 para variante secuencial) \\
\texttt{road\_length}   & Tamaño de la carretera $N$ \\
\texttt{density}        & Densidad inicial de carros $\rho$ \\
\texttt{repetition}     & Índice de repetición (1 a 10) \\
\texttt{wall\_time\_ms} & Tiempo máximo entre procesos, ecuación~\eqref{eq:tiempo_mpi} \\
\texttt{avg\_velocity}  & Velocidad media $\bar{v}$ sobre los $T_m$ pasos \\
\hline
\end{tabular}
\end{center}

El tiempo reportado corresponde al proceso más lento conforme a la
ecuación~\eqref{eq:tiempo_mpi}, lo que refleja el tiempo real de cada
ejecución sincronizada y es el valor apropiado para calcular el speedup
respecto a la línea de referencia secuencial.